
\textbf{Interpretation.} The results validate the core insight that MCP tool-calling traces encode researcher reasoning in natural language, enabling zero-annotation semantic extraction for two of three layers (syntactic validation and semantic content extraction). However, the complete failure of constraint inference (Layer 3) reveals a critical gap between semantic extraction and automated contradiction detection.

The 97.48\% natural language presence rate in both query and result (exceeding the 80-90\% estimate) likely reflects research domain bias rather than a universal MCP property. Research pipelines are inherently NL-rich---queries are hypothesis-driven, results are findings. Production pipelines with structured data processing may exhibit lower NL content (60-70\%), limiting generalization. This scopes the findings: the validated extraction quality (82.7\%/86.3\%) applies to research pipelines, not universally to all MCP use cases.

The 82.7\% recall and 86.3\% precision achieved with simple few-shot prompts (no fine-tuning, no labeled training data) suggests pre-trained LLM capability on scientific text. Claude Sonnet 4.5's training corpus includes academic papers, technical documentation, and research artifacts, making it well-suited for assumption/claim extraction from research traces.

The 0\% recall across 1,200 assumption-claim pairs aligns with known limitations in NLP literature. Cosine similarity on sentence embeddings optimizes for semantic relatedness (topic clustering, paraphrase detection), not logical contradiction. The statements ''effective rank will decrease'' and ''effective rank increased by 6.02\%'' have high cosine similarity (>0.3) because they share the same entity (effective rank) and concept (change). Detecting that they contradict requires entailment models (BERT-NLI, DeBERTa-MNLI) that explicitly model logical relationships (entailment, contradiction, neutrality). This negative result extends findings on Natural Language Inference---embedding-based similarity is insufficient for NLI tasks.

\textbf{Limitations.} The constraint inference method failure (0\% recall) indicates that the semantic similarity approach is fundamentally unsuited for contradiction detection, as cosine distance measures topic relatedness rather than logical polarity. While this is a failure of the current method, it is a valuable negative result---it establishes a clear boundary condition (embeddings cannot detect contradictions) and points toward entailment models or LLM-based reasoning as necessary alternatives. The failure is methodological, not a flaw in the broader hypothesis that MCP traces encode testable constraints.

Test data scope is limited to research pipelines. All 20 MCP traces come from hypothesis generation, literature search, and experiment design workflows. The 97.48\% NL presence rate may not generalize to production MCP use cases (automation scripts, data pipelines) that use structured tool calls with minimal descriptive text. This limitation is acceptable because the hypothesis explicitly scoped to research pipelines. If applying to production pipelines, validation should be re-run on domain-specific traces.

The constraint inference validation used only one documented contradiction (h-m1 effective rank case) as ground truth (N=1). While 0\% recall on N=1 is statistically weak evidence, the threshold tuning analysis compensates---testing 1,200 pairs across five thresholds (0.2-0.4) showed zero detections even at the loosest setting. This suggests the issue is systematic method mismatch, not insufficient ground truth.

End-to-end validation was not performed because Layer 3 constraint inference failed its gate. The full three-layer framework's predicted failure detection rate cannot be verified without functional constraint inference. This dependency-driven blocking is intentional---if Layer 3 fails, end-to-end testing would produce misleading results. The pipeline correctly stopped forward progress rather than compounding errors.

\textbf{Connection to Existing Literature.} The finding that MCP traces encode rich NL content (97.48\% presence) extends Ahn et al. (2025)'s work on MCP for medical concept standardization. While Ahn used MCP for tool composition in a structured domain, this work validates traces themselves as semantic artifacts suitable for validation-as-inference, opening a new use case for MCP beyond tool orchestration.

Zero-annotation extraction (82.7\%/86.3\%) complements Fu et al. (2025)'s PRDBench. Fu reduces annotation cost via agent-driven benchmark generation; this work achieves zero-annotation by extracting from execution traces. Both leverage LLM capability but attack different points in the validation pipeline---Fu generates test cases, this work infers constraints from runtime behavior.

The constraint inference failure (0\% recall with semantic similarity) aligns with NLI literature. Bowman et al. (2015) SNLI and Williams et al. (2018) MultiNLI demonstrate that embedding-based similarity is insufficient for entailment and contradiction tasks, which require models explicitly trained on logical relationships. The negative result empirically confirms this limitation in the research pipeline domain.

The approach diverges from Neutatz et al. (2021)'s declarative constraints for ML. Neutatz requires manual declaration of feature constraints, achieving reliable enforcement via explicit verification. This work infers constraints from traces (no manual declaration), but Layer 3 failed---suggesting a hybrid approach (declarative + inferred) may be necessary for production use.

\textbf{Implications.} The validated Layers 1-2 (syntactic validation and semantic extraction) can be deployed immediately for practical use: trace quality checks (completeness monitoring), assumption extraction (hypothesis archaeology from logs), and claim extraction (results summarization). These components achieve human-competitive quality (82.7\%/86.3\%) without manual annotation.

Layer 3 (constraint inference) requires redesign with entailment models (BERT-NLI, DeBERTa-MNLI) or LLM-based contradiction detection. The path forward is clear: replace sentence-transformer similarity with probability of contradiction from an NLI model, using a threshold to flag mismatches. This addresses the root cause (semantic similarity measures topic relatedness, not logical polarity) without requiring changes to Layers 1-2 or the broader framework architecture.
